\section{Introduction}
\label{sec:intro}

Five labelled aerial scenes of a class the network has never seen. A ResNet-18 trained on $25$ other NWPU-RESISC45 classes labels the query scenes of the new class correctly \num{resisc45:dense_ce:novel_5}\% of the time. Keep $30\%$ of its parameters and the figure falls to \num{resisc45:A_dacis_PMP_30:novel_5}\% with the DACIS score, \num{resisc45:A_random_PMP_30:novel_5}\% with randomly chosen channels and \num{resisc45:A_taylor_PMP_30:novel_5}\% with a first-order Taylor score. This paper asks whether such differences between channel scores are real, and what else moves the number.

The question matters where a classifier has to run on small hardware that cannot be retrained often: the payload of a small satellite, a drone over a field. On-board inference already flies~\cite{giuffrida2020cloudscout,giuffrida2022phisat}, the classes that matter change with the mission, and labels for a new class are few. A pruned network is cheap to run and a prototype head can be updated from five labelled patches without touching the weights. We measure cost only on a server CPU and GPU, so this motivation is context, not a result. The pruning studies we build on~\cite{li2017pruning,liu2017slimming,molchanov2017pruning,he2017channel} report accuracy on the classes the network was trained for. What survives when the classes at test time are not those the channels were chosen for?

We study that on ResNet-18. Every method removes the inner convolution channels of each residual block under one shared depth profile (Sec.~\ref{sec:method}), so two scores differ only in which channels they keep. The scores are random selection, filter $\ell_1$ norm, BatchNorm scale, a Taylor term, and DACIS, a rank sum of gradient norm, activation variance and Fisher ratio over class labels, which we also run in its three parts. DACIS and the PMP schedule of prune, meta-train, prune again come from an earlier preprint~\cite{uddin2026preprint}; this paper re-tests them with held-out classes, five class partitions and shared test episodes, and they do not hold up. We test on held-out classes of aerial scenes (NWPU-RESISC45~\cite{cheng2017resisc}) and of Sentinel-2 patches with all thirteen bands (EuroSAT-MS~\cite{helber2019eurosat}). PlantVillage leaf photographs~\cite{hughes2015plantvillage}, with PlantDoc field photographs~\cite{singh2020plantdoc} as a shift check, serve as a ground-level control because the earlier preprint used them; they lack geospatial context and fall outside the Earth-observation scope of the workshop.

The contributions are four findings and the protocol that supports them.
\begin{itemize}[leftmargin=*,itemsep=0pt,topsep=2pt]
  \item Under PMP at $30\%$ of the parameters, \num{cmp:above:PMP} of \num{cmp:total:PMP} paired comparisons of a computed score with random selection favour the computed score and \num{cmp:below:PMP} favour random selection. Only the Taylor term is ahead, by at most \num{cmp:maxgain:PMP} points, and no comparison survives a Bonferroni correction. DACIS cannot be told from random selection.
  \item Episodic training does not raise the accuracy of the dense network, and the small gain of two-stage schedules is matched by a variant that updates no weight in the meta stage.
  \item Pruning saves less time than operations: \num{cost:macs-red:30}\% fewer multiply-accumulates gave \num{cost:cpu-red:30}\% lower single-thread CPU latency. The budget moves accuracy several times more than the score.
  \item An unpruned MobileNetV3-Small, trained with the same recipe, beats the pruned and the dense ResNet-18 on both Earth-observation sets at a fraction of the cost. This is a strong baseline that pruning studies in this setting should include.
\end{itemize}
The protocol (class-disjoint partitions, shared test episodes, paired intervals over partitions, the count of significant comparisons set against what chance alone gives) is meant to be reused. With five partitions it cannot exclude differences of about \num{cmp:halfw:PMP} points, and Sec.~\ref{sec:limits} lists what we did not test.
